\subsection{A pendant singleton at a small tree corner}

\begin{proposition}[Proved: a small-corner pendant in the eighteen-arrow budget]
\label{CP2:A-pendant-small-tree-eighteen}
Let $C$ be a finite connected category with at least three
objects, acyclic support, and at most eighteen arrows.
Suppose its undirected comparability graph is a tree.
Assume that a leaf $p$ has trivial corner and a singleton
stalk to its neighbor $v$, whose corner $N$ has order at
most four. Suppose that, apart from $v$, at most one corner
has order greater than two, and that corner has order at
most eleven. Then $BC$ does not have type $\CP^2$.
\end{proposition}
\begin{proof}
The full edge on $p,v$ has initial or terminal object $p$,
according to its orientation, so its classifying space is
contractible. Let $D=C\setminus\{p\}$ be the full deletion.
There is no strict three-object chain in a comparability
tree. Hence the nerves of $D$ and that edge cover $NC$,
and their intersection is $NN$.

If $BN$ is contractible, this is a homotopy pushout along
CW cofibrations, and the actual inclusion $BD\to BC$ is
a homotopy equivalence. The deletion removes exactly the
identity of $p$ and its sole cross arrow, so $D$ has at
most sixteen arrows. Theorem~\ref{CP2:A-global-lower-seventeen}
excludes projective-plane type.

Otherwise $BN$ has the homotopy type of a nontrivial finite
group by Corollary~\ref{CP2:A-five-types}. The minimal ideal
of $N$ is a group: two rows or two columns with a nontrivial
local group would require at least four ideal elements,
and the outside identity would give order at least five.
Both trivial rational modules of $N$ are projective by
averaging over this group ideal.

All edges at $v$ have the same orientation, since any
incoming and outgoing pair would create a transitive
triangle. Apply the appropriately sided actual orbit
functor of Lemma~\ref{CP2:A-high-orbits}, replacing $N$
by the trivial corner and all its stalks by their orbit
sets. Its pullback
$H^*(B\overline C;\mathbb Q)\to H^*(BC;\mathbb Q)$
is a ring isomorphism. Orbit sets of nonempty stalks
are nonempty, so the comparability tree is unchanged.
All corners of $\overline C$, except possibly one of
order at most eleven, have order at most two. Under the
assumption $BC\simeq\CP^2$, rational homology of
$B\overline C$ is bounded. Proposition~\ref{CP2:A-tree-one-large-corner-eleven}
therefore kills all degree-two products there and, by
pullback, on $BC$. This contradicts the primitive rational
square of the projective plane.
\end{proof}
